\subsection{墙与面}

定义 3. 设 C 是 E 的一个室。C 的一个面是包含于 C 的闭包且支撑为超平面的面片。C 的一面墙是某个 C 的面的支撑超平面。

C 的每一面墙都属于 H。命题 4 表明，超平面 \(L \in H\) 是 C 的一面墙，当且仅当 \(C \neq D_{\mathfrak{H} - \{L\}}(C)\)。此外，C 的每一面墙都是 C 的唯一一个面的支撑。

命题 8. 属于 H 的每个超平面 H 至少是一个室的一面墙。

由命题 7 (iii)，\(a\) 中存在一点 \(\mathbf{H}\)，不属于 \(\mathbf{H}' \neq \mathbf{H}\) 中任何超平面 \(\mathfrak{H}\)；由命题 6，存在一个室 \(\mathbf{C}\)，使得 \(a \in \overline{\mathbf{C}}\)；于是命题 4 表明 \(\mathbf{H}\) 是 \(\mathbf{C}\) 的一面墙。

命题 9. 设 C 是一个室，M 是 C 的墙集。则 \(\mathrm{C} = \mathrm{D}_{\mathfrak{M}}(\mathrm{C})\)，并且 H 的每个满足 \(\mathrm{C} = \mathrm{D}_{\mathfrak{L}}(\mathrm{C})\) 的子集 L 都包含 M。\(\overline{C}\) 的子集 F 是面片，当且仅当它是相对于族 M 的 E 的面片。

\begin{enumerate}
\def\labelenumi{\alph{enumi})}
\item
  设 \(\mathfrak{L}\) 是 \(\mathfrak{H}\) 的一个子集，满足 \(C = D_{\mathfrak{L}}(C)\)。考虑属于 \(L\) 但不属于 \(\mathfrak{H}\) 的超平面 \(\mathfrak{L}\)；令 \(\mathfrak{N}\) 为属于 \(H \neq L\) 的超平面 \(\mathfrak{H}\) 所成的集合。则 \(\mathfrak{L} \subset \mathfrak{N}\)，从而 \(C = D_{\mathfrak{M}}(C)\)，且 \(L\) 不与 \(D_{\mathfrak{M}}(C)\) 相交。由命题 4 中蕴含 \(\Longrightarrow\)，超平面 \(L\) 不是 \(C\) 的墙。因此，\(C\) 的每一面墙都属于 \(\mathfrak{L}\)。
\item
  假设 \(C = D_{\mathfrak{L}}(C)\)。设 H 是属于 L 但不是 C 的墙的超平面，置 \(L' = L - \{H\}\)。由命题 4 中蕴含 (iii) \(\Longrightarrow\) (i)，凸集 \(D_{\mathfrak{L}'}(C)\) 不与 H 相交，所以 \(D_{\mathfrak{L}'}(C) \subset D_{H}(C)\)，且 \(C = D_{\mathfrak{L}'}(C)\)。若 F 是 L 的不包含 C 的任何墙的有限子集，则对 F 的基数作归纳，得到 \(C = D_{\mathfrak{L} - \mathfrak{F}}(C)\)。
\item
  设 \(a\) 是 C 的一点；显然，\(\mathrm{C} \subset \mathrm{D}_{\mathfrak{M}}(a)\)。设 \(a'\) 是 \(\mathrm{D}_{\mathfrak{M}}(a)\) 的一点；由于闭线段 \([aa']\) 是紧的，与 \(\mathfrak{F}\) 相交的超平面 \(\mathrm{H} \in \mathfrak{H}\) 所成的集合 \([aa']\) 是有限的。由于 \(a\) 和 \(a'\) 严格位于 C 的每一面墙的同侧，C 的墙没有一个属于 \(\mathfrak{F}\)；由 \(b\)，得 \(\mathrm{C} = \mathrm{D}_{\mathfrak{H}-\mathfrak{F}}(\mathrm{C})\)。由于 \(a' \in \mathrm{D}_{\mathfrak{H}-\mathfrak{F}}(a)\)，有 \(a' \in \mathrm{C}\)。因此我们证明了 \(\mathrm{D}_{\mathfrak{M}}(a) \subset \mathrm{C}\)，这就证明了命题的第一部分。
\item
  为证明命题的最后一个断言，显然只需证明：\(\overline{C}\) 的一个相对于 M 为 E 的面片的子集 F，也是一个相对于 H 为 E 的面片；或者说，与 F 相交的每个超平面 \(H \in H\) 都包含 F。于是设 H 是一个与 F 相交但不包含 F 的超平面。由于 F 在其仿射支撑中是开集，它不完全位于 H 的一侧。由此 \(\overline{C}\) 也不完全位于 H 的一侧，因而超平面 H 不属于 H，证明完成。
\end{enumerate}

注. 1) 由公式 (6) 和命题 9 可知，一个室 C 的闭包是由 C 的墙界定且包含 C 的闭半空间之交。

\begin{enumerate}
\def\labelenumi{\arabic{enumi})}
\setcounter{enumi}{1}
\tightlist
\item
  设 \(\mathbf{F}\) 是一个支撑为超平面 \(\mathbf{L}\) 的面片；我们将证明 \(\mathbf{F}\) 是两个室的面。令 \(\mathfrak{N}\) 为属于 \(\mathrm{H} \neq \mathrm{L}\) 的超平面 \(\mathfrak{H}\) 所成的集合；置 \(\mathrm{A} = \mathrm{D}_{\mathfrak{N}}(\mathrm{F})\)，并用 \(\mathrm{D}^{+}\) 和 \(\mathrm{D}^{-}\) 表示由 \(\mathrm{L}\) 界定的两个开半空间。集合 \(\mathbf{A}\) 是开集并包含 \(\mathbf{F} \subset \mathbf{L}\)，而由于 \(\mathbf{L}\) 的每一点都属于 \(\mathrm{D}^{+}\) 和 \(\mathrm{D}^{-}\) 的闭包，集合 \(\mathrm{C}^{+} = \mathrm{A} \cap \mathrm{D}^{+}\) 和 \(\mathrm{C}^{-} = \mathrm{A} \cap \mathrm{D}^{-}\) 非空；它们是室。此外，超平面 \(\mathbf{L}\) 与 \(\mathrm{D}_{\mathfrak{N}}(\mathrm{F}) = \mathrm{D}_{\mathfrak{N}}(\mathrm{C}^{+})\) 相交；命题 4 表明 \(\mathbf{L}\) 是 \(\mathrm{C}^{+}\) 的一面墙，并且
\end{enumerate}

F 与 \(L \cap D_{\mathfrak{N}}(F)\) 相交，是 \(C^{+}\) 的一个面；同样，F 是 \(C^{-}\) 的一个面。最后，设 C 是以 F 为面的一个室，并且例如假设 \(D^{+} = D_{L}(C)\)；由命题 4，集合 \(D_{\mathfrak{N}}(C)\) 与 F 相交，因而等于 \(D_{\mathfrak{N}}(F)\)，于是有

\[
C = D _ {\mathfrak {H}} (C) = D _ {L} (C) \cap D _ {\mathfrak {N}} (C) = D ^ {+} \cap D _ {\mathfrak {N}} (F) = C ^ {+}.
\]

